\documentclass{beamer}


\usepackage{graphicx}
\usepackage{pgffor}
\usepackage{tikz}
\usepackage{geometry}

% The original pages are 792 x 612 bp.  Removing 38 bp from the bottom
% gives a 792 x 574 bp slide with no navigation-symbol footer.
% \geometry{paperwidth=792bp,paperheight=574bp}
\setbeamertemplate{navigation symbols}{}
\setbeamertemplate{footline}{}

% Keep this PDF in the same directory as the .tex file.
\newcommand{\sourcepdf}{Source_Lecture01_pt2/Lecture01_pt2_footer_removed.pdf}
\newcommand{\bottomcrop}{30bp}

\usepackage[utf8]{inputenc}
\usepackage{amsmath}
\usepackage{enumitem}
\usepackage{ragged2e}
\usepackage{listings}
\usepackage{minted}
\usepackage{animate}
\AtBeginEnvironment{frame}{\justifying}
\AtBeginEnvironment{itemize}{\justifying}
\AtBeginEnvironment{enumerate}{\justifying}
\usepackage{tcolorbox}
\tcbuselibrary{listings, minted}

\setminted{
  fontsize=\small,
  linenos,
  breaklines
}
\usetheme{Madrid}
\usepackage{amsmath, amsfonts, amssymb}

\setbeamertemplate{footline}[frame number]
\setbeamertemplate{navigation symbols}{}

\usetikzlibrary{positioning}

\definecolor{ai}{RGB}{14,79,132}
\definecolor{ml}{RGB}{38,126,186}
\definecolor{dl}{RGB}{86,176,232}
\usetikzlibrary{calc}
\usetikzlibrary{arrows.meta}

\definecolor{lightblue}{RGB}{200,220,235}

\renewcommand{\familydefault}{\sfdefault}

\usetikzlibrary{arrows.meta}

\definecolor{classicalyellow}{RGB}{249,227,150}
\definecolor{machineblue}{RGB}{205,224,241}

\title[CSC-220: Neural Networks and Deep Learning]{CSC-220: Neural Networks and Deep Learning}
\subtitle{Lecture 01 Part II - DL: Mathematical Building Blocks, Gradient Descent, and Backpropagation}

\date[Sep. 15\textsuperscript{th}, 2026]{September 15\textsuperscript{th}, 2026}

\setbeamertemplate{footline}{
  \leavevmode%
  \hbox{%
    \begin{beamercolorbox}[wd=.4\paperwidth,ht=2.5ex,dp=1ex,left]{author in head/foot}%
      \hspace{1em}\insertshorttitle
    \end{beamercolorbox}%
    \begin{beamercolorbox}[wd=.4\paperwidth,ht=2.5ex,dp=1ex,center]{date in head/foot}%
      \insertshortdate
    \end{beamercolorbox}%
    \begin{beamercolorbox}[wd=.2\paperwidth,ht=2.5ex,dp=1ex,right]{date in head/foot}%
      \insertframenumber{} / \inserttotalframenumber\hspace{1em}
    \end{beamercolorbox}%
  }%
}

\begin{document}

\thispagestyle{empty}
\frame{\titlepage}

\setcounter{framenumber}{0}

\begin{frame}{How Learning (Backpropagation) Works}
    \centering
    \animategraphics[
        controls,
        loop,
        poster=first,
        width=.86\paperwidth
    ]{3}{backprop_frames/backprop-}{000}{069}
\end{frame}

\begin{frame}{Let's zoom in \dots}
    \centering
    \animategraphics[
        controls,
        loop,
        poster=first,
        width=.86\paperwidth
    ]{4}{backprop_gradient_animation/backprop_gradient_frames/gradient-}{000}{083}
\end{frame}



\foreach \pdfpage in {1,...,5}{%
        \begin{frame}
            \begin{tikzpicture}[remember picture,overlay]
                \node[inner sep=0pt] at (current page.center) {%
                    \includegraphics[
                        page=\pdfpage,
                        % trim=0 \bottomcrop{} 0 0,
                        clip,
                        width=\paperwidth
                    ]{\sourcepdf}%
                };
            \end{tikzpicture}
        \end{frame}
    }

\begin{frame}{ReLU functions and backpropagation}
    \begin{block}{Key idea}
        For a pre-activation \(a=\mathbf{w}^{\top}\mathbf{x}+b\),
        \[
            \operatorname{ReLU}(a)=\max(0,a).
        \]
        Positive signals and their gradients pass through the neuron.
        Negative signals are mapped to zero and block the gradient for
        that input.
    \end{block}
\end{frame}

\begin{frame}{ReLU functions and backpropagation}
    \centering
    \animategraphics[
        controls,
        loop,
        poster=first,
        width=.94\paperwidth
    ]{12}{relu_bp/gradient-}{000}{083}
\end{frame}

\foreach \pdfpage in {7,...,57}{%
        \begin{frame}
            \begin{tikzpicture}[remember picture,overlay]
                \node[inner sep=0pt] at (current page.center) {%
                    \includegraphics[
                        page=\pdfpage,
                        % trim=0 \bottomcrop{} 0 0,
                        clip,
                        width=\paperwidth
                    ]{\sourcepdf}%
                };
            \end{tikzpicture}
        \end{frame}
    }

\end{document}
